\DocumentMetadata{
  pdfversion=2.0,
  pdfstandard=ua-2,
  lang=en-US,
  tagging=on,
  tagging-setup={math/setup=mathml-SE,tabsorder=structure}
}
\documentclass[11pt,letterpaper]{article}
\usepackage[margin=0.68in,includefoot]{geometry}
\usepackage{newcomputermodern}
\usepackage{amsmath,amscd,mathtools}
\usepackage{tikz,adjustbox}
\providecommand{\square}{\mdlgwhtsquare}
\usepackage[hidelinks]{hyperref}
\hypersetup{
  pdftitle={Topology PhD Exam, May 1993},
  pdfauthor={University of Florida Department of Mathematics},
  pdfsubject={Graduate mathematics examination},
  pdfdisplaydoctitle=true,
  bookmarksopen=true
}
\setcounter{secnumdepth}{0}
\setlength{\parindent}{0pt}
\setlength{\parskip}{0.28em}
\setlength{\emergencystretch}{3em}
\setlength{\leftmargini}{2.15em}
\setlength{\leftmarginii}{2.65em}
\setlength{\leftmarginiii}{3em}
\renewcommand{\labelenumi}{\textbf{\theenumi.}}
\pagestyle{empty}
\raggedbottom
\allowdisplaybreaks
\newenvironment{examproblems}[1]{%
  \begin{enumerate}%
  \setcounter{enumi}{#1}%
  \setlength{\topsep}{0.35em}%
  \setlength{\partopsep}{0pt}%
  \setlength{\itemsep}{0.48em}%
  \setlength{\parsep}{0.18em}%
}{\end{enumerate}}
\newenvironment{examsubparts}{%
  \begin{enumerate}%
  \setlength{\topsep}{0.18em}%
  \setlength{\partopsep}{0pt}%
  \setlength{\itemsep}{0.16em}%
  \setlength{\parsep}{0.1em}%
}{\end{enumerate}}
\newenvironment{examinstructions}{%
  \par\begingroup\itshape
}{%
  \par\endgroup\vspace{0.18em}
}
\makeatletter
\renewcommand\section{\@startsection{section}{1}{0pt}%
  {-1.25ex plus -.25ex minus -.1ex}%
  {0.55ex plus .1ex}%
  {\normalfont\large\bfseries}}
\renewcommand{\@maketitle}{%
  \newpage\null\vskip -1em%
  {\footnotesize\bfseries
    \href{https://www.ufl.edu/}{University of Florida}, \href{https://math.ufl.edu/}{Department of Mathematics}\par}%
  \vskip -0.5em%
  \tagstructbegin{tag=Title}%
    {\LARGE\bfseries\href{https://gma.math.ufl.edu/exams/topology-phd/}{Topology PhD Exam}, May 1993\par}%
  \tagstructend%
  \vskip 0.25em%
  \tagmcbegin{artifact}%
    \hrule height 1pt%
  \tagmcend%
  \vskip 1em%
}
\makeatother
\title{Topology PhD Exam, May 1993}
\author{}
\date{}
\begin{document}
\maketitle
\thispagestyle{empty}
\section{PART I}
\begin{examinstructions}
State the following theorems:
\end{examinstructions}
\begin{examproblems}{0}
\item
Baire Category Theorem
\item
Urysohn’s Lemma
\item
Brouwer Fixed Point Theorem
\item
Jordan Separation Theorem for \(R^2\)
\item
Schoenflies Theorem for \(R^2\). Does the corresponding statement hold for \(R^3\)? Explain.
\end{examproblems}
\section{PART II}
\begin{examinstructions}
For the following problems, show all work and write clearly. The proofs should be detailed, leaving no doubt about the correctness!
\end{examinstructions}
\begin{examproblems}{0}
\item
What is the universal cover of \(P^2\)? Prove your answer.
\item
Prove: Let~\(C\) be the Cantor set, and let~\(X\) be any compact metric space. Then there exists a continuous, surjective function \(f:C\to X\).
\item
State and prove the Finite Simplicial Approximation Theorem.
\item
Use the Seifert–van Kampen Theorem to compute the fundamental group of the Klein bottle.
\item
This problem has two parts.
\begin{examsubparts}
\item[\textnormal{(i)}]
What is \(\widetilde H_0(X)\), when~\(X\) is a two point space? (Do not prove.)
\item[\textnormal{(ii)}]
Use the Mayer–Vietoris sequence and (i) to compute the homology of \(S^n\), for \(n\ge 1\).
\end{examsubparts}
\end{examproblems}
\end{document}
